\section*{Límites y Alcances}
\addcontentsline{toc}{section}{Límites y Alcances}

El sistema abarca el recorrido completo desde una descripción en lenguaje natural hasta un netlist simulado en ngspice. Del lado del cliente, gestiona el registro y el inicio de sesión, los espacios de trabajo y las claves con las que el usuario envía sus solicitudes. Del lado del servidor, lee la descripción y extrae de ella los parámetros eléctricos, divide los circuitos de varias etapas en sub-bloques, calcula los valores de los componentes, genera el código PySpice y el netlist, simula en ngspice, lee las métricas, ajusta los valores de forma iterativa a través del curador y entrega al usuario el netlist final junto con sus métricas y el historial de iteraciones.

El sistema trabaja únicamente con circuitos analógicos CA/CD formados por componentes comerciales que admiten un modelo SPICE, como resistores, capacitores, inductores, fuentes, diodos, transistores y amplificadores operacionales descritos por su macromodelo. Un componente sin un modelo utilizable queda fuera de su alcance, ya que sin modelo no hay simulación y, por tanto, no hay señal de realimentación. El sistema no diseña la disposición física de la placa (PCB), no realiza síntesis a nivel de transistor para un proceso de fabricación específico (PDK) ni genera esquemáticos gráficos. Por último, los valores numéricos que aparecen en los requisitos son metas de la primera versión y pueden ajustarse durante la evaluación experimental.
